\chapter{TikZ Drawing Guide}
\label{chap:tikz-guide}

Every diagram in this book is built from a small set of reusable patterns.
This appendix collects them so that figures added for an exercise or a new
section stay visually consistent with the book.  The examples below are
adapted from figures already used in the chapters; the key ideas are the
\emph{state box}, the \emph{arrow with a label}, the \emph{pipeline row},
and the \emph{timeline}.

\section{The Shared Preamble and Styles}

The book's preamble loads the libraries that all figures use
(\texttt{positioning}, \texttt{arrows.meta}, \texttt{calc},
\texttt{shapes.geometric}).  Reuse the style names defined here so that
figures share one look:

\begin{lstlisting}[language={}, caption={Common TikZ styles used throughout the book.}]
\begin{tikzpicture}[
    node distance=1.8cm and 2.2cm,
    state/.style={draw, rounded corners, align=center,
                  minimum width=2.6cm, minimum height=0.9cm},
    edge/.style={-{Latex}, thick},
    abort/.style={-{Latex}, thick, red}
]
\end{tikzpicture}
\end{lstlisting}

\emph{Style guide:} draw states in black, normal transitions with
\texttt{-{Latex}} in black, and \emph{abort} or \emph{conflict} edges in
red (\texttt{abort/.style}).  Reserve a third accent colour (e.g. blue) for
a second, parallel path such as the fallback path in a hybrid design
(Figure~\ref{fig:phased-hybrid-state}).  Keep fonts at
\texttt{\textbackslash footnotesize} or smaller on edge labels.

\section{Pattern 1: A Pipeline Row}

A pipeline is a row of boxes with labelled arrows.  See
Figure~\ref{fig:front-end-pipeline} (the front end) and
Figure~\ref{fig:compiler-tm-pass} (the compiler pass) for the pattern in
use:

\begin{lstlisting}[language={}, caption={Pipeline row.}]
\begin{tikzpicture}[box/.style={draw, rounded corners, align=center,
    minimum width=2.2cm, minimum height=0.8cm}]
\node[box] (src) {Source};
\node[box, right=of src] (lex) {Lexer};
\node[box, right=of lex] (ast) {AST};
\draw[-{Stealth}] (src) -- node[above]{tokens} (lex);
\draw[-{Stealth}] (lex) -- node[above]{parse} (ast);
\end{tikzpicture}
\end{lstlisting}

To turn a row into a ring (the retry loop of an abort), route the last box
back under the row with \texttt{(abort.west) -| (begin.south)} as in
Figure~\ref{fig:htm-coherence-abort}.

\section{Pattern 2: A State Machine}

A finite-state machine is the book's most common diagram: the transaction
lifecycle (Figure~\ref{fig:transactional-state-machine}), the API control
states (Figure~\ref{fig:api-state-machine}), and the resource state machine
(Figure~\ref{fig:tx-resource-state}) all use it.

\begin{lstlisting}[language={}, caption={State machine with a self-loop and a labelled decision.}]
\begin{tikzpicture}[
    node distance=2cm,
    state/.style={draw, rounded corners, align=center,
                  minimum width=2.2cm, minimum height=0.8cm},
    edge/.style={-{Latex}, thick}
]
\node[state] (begin) {Begin};
\node[state, right=of begin] (read) {Read};
\node[state, right=of read] (commit) {Commit};
\node[state, below=of read] (abort) {Abort};
\draw[edge] (begin) -- (read);
\draw[edge] (read) -- node[above]{valid} (commit);
\draw[edge] (read) -- node[right]{stale} (abort);
\draw[edge] (abort.west) -| ($(begin.south)+(0,-0.4)$) -- (begin.south);
\end{tikzpicture}
\end{lstlisting}

For a Markov chain (Chapter~5, \S\ref{sec:math-example}), put the
transition probability on every edge and keep the chain acyclic except for
the absorbing commit/abort states.

\section{Pattern 3: A Timeline}

A timeline shows when reads observe which clock values, or when a commit
publishes writes.  It is the pattern behind the TL2 timeline
(Figure~\ref{fig:tl2-timeline}) and the read-validate diagrams.

\begin{lstlisting}[language={}, caption={A timeline with read and commit markers.}]
\begin{tikzpicture}[x=1.4cm]
\draw[-{Latex}] (0,0) -- (5,0) node[right]{time};
\foreach \t/\v in {1/1,2/2,3/3} {
    \draw (\t,0.08) -- (\t,-0.08) node[below, font=\scriptsize]{clock=\v};
}
\node[above, font=\small] at (1.5,0.1) {read};
\node[above, font=\small] at (3,0.1) {validate};
\end{tikzpicture}
\end{lstlisting}

\section{Pattern 4: An Abort/Retry Waveform}

The waveform that shows an abort, a retry, and a later commit is a row of
short horizontal segments with a drop at the abort point.  Build it from a
chain of \texttt{edge} arrows where the \emph{abort} edge uses the red
style and the retry edge arcs back, as in
Figure~\ref{fig:phased-hybrid-state}.

\section{Pattern 5: Warp Divergence and Ballot (Part VII)}

GPU diagrams (Chapter~17) draw a warp as a column of lanes, with a
\emph{divergence} branch shown as two forks.  Keep lanes aligned with
\texttt{below=of} and mark inactive lanes with a dashed outline.  For a
ballot result, show the per-lane predicate bits feeding a mask, as in
\S\ref{sec:ballot}.

\section{Adding a New Figure}

\begin{enumerate}
    \item Pick the pattern above that matches the concept (pipeline,
          state machine, timeline, or waveform).
    \item Copy its template and rename every label.
    \item Add a \texttt{\textbackslash caption} that states the one
          conclusion the reader should take away, and a
          \texttt{\textbackslash label} in the figure's convention
          (\texttt{fig:}).
    \item Reuse the shared styles so the figure matches the book's palette.
\end{enumerate}

The figures already in the book are the reference implementation of these
patterns.  When in doubt, copy the nearest existing figure and modify it
rather than inventing new styling.